\section{Encoder Architecture}
\label{sec:architecture}

\subsection{Optimisation Objective}

Given a greyscale target image $\I \in [0,1]^{H \times W}$, the encoder finds
physical Gaussian parameters $\W^* \in \R^{K \times 7}$ that minimise
mean-squared error (MSE) between the rendered image and the target:

\begin{equation}
  \W^* = \argmin_{\W} \;\MSE\!\bigl(\Render(\W),\, \I\bigr)
       = \argmin_{\W} \;\frac{1}{HW}\sum_{i,j}
         \bigl(\Render(\W)_{ij} - \I_{ij}\bigr)^2,
  \label{eq:objective}
\end{equation}

\noindent where $\Render(\W)$ is the differentiable Gaussian splatting renderer
(Section~\ref{sec:renderer}).

\subsection{Parameter Layout}

Each Gaussian $k \in \{1,\ldots,K\}$ is described by a 7-dimensional parameter
vector $\W_{k,:} = [\sigma_x,\, \sigma_y,\, \rho,\, \alpha,\, c,\, x,\, y]$
with the physical ranges shown in Table~\ref{tab:params}.

\begin{table}[ht]
\centering
\caption{Physical parameters of each 2-D Gaussian (one row of $\W$).}
\label{tab:params}
\begin{tabular}{lllp{6.5cm}}
\toprule
\textbf{Symbol} & \textbf{Name} & \textbf{Range} & \textbf{Meaning} \\
\midrule
$\sigma_x$ & std.\ dev.\ $x$ & $(0, 1)$  & Horizontal spread of the Gaussian \\
$\sigma_y$ & std.\ dev.\ $y$ & $(0, 1)$  & Vertical spread of the Gaussian \\
$\rho$     & correlation     & $(-1, 1)$ & Off-axis tilt; covariance $= \rho\sigma_x\sigma_y$ \\
$\alpha$   & opacity         & $(0, 1)$  & \emph{Vestigial} — not passed to renderer (§\ref{sec:alpha}) \\
$c$        & colour          & $(0, 1)$  & Greyscale intensity contributed by this Gaussian \\
$x$        & position $x$    & $(-1, 1)$ & Horizontal centre in normalised image coords \\
$y$        & position $y$    & $(-1, 1)$ & Vertical centre in normalised image coords \\
\bottomrule
\end{tabular}
\end{table}

\subsection{Raw Parameterisation and Differentiable Transforms}

The optimiser operates on \emph{unconstrained raw parameters}
$\Wraw \in \R^{K \times 7}$, mapped to physical space by smooth, bounded
transforms that maintain differentiability throughout:

\begin{align}
  \sigma_x &= \sigmoid(\Wraw_{:,0}), &
  \sigma_y &= \sigmoid(\Wraw_{:,1}), \label{eq:sigma}\\
  \rho     &= \tanh(\Wraw_{:,2}),    &
  \alpha   &= \sigmoid(\Wraw_{:,3}), \label{eq:rho}\\
  c        &= \sigmoid(\Wraw_{:,4}), &
  x        &= \tanh(\Wraw_{:,5}),\quad
  y = \tanh(\Wraw_{:,6}).              \label{eq:xy}
\end{align}

\noindent This reparameterisation ensures that gradient descent on $\Wraw$
cannot push any parameter outside its valid range, eliminating the need for
projected gradient methods or explicit clipping during training.

\subsection{Initialisation}
\label{sec:init}

Initialisation has a large impact on convergence speed.
Gradient descent on the rendering objective has many local minima, and
Gaussians that start far from bright image regions converge slowly or not at all.

\paragraph{Brightness-weighted seeding.}
We seed position $(x_k, y_k)$ by sampling pixel coordinates from the
distribution $p_\text{pixel} \propto \I_{ij}$ with replacement.
This concentrates initial Gaussians in bright (ink) regions where they must
contribute, rather than in background regions where their gradient signal is
near-zero.

\medskip\noindent
\textbf{Why with replacement?}\ \ Sampling without replacement spreads Gaussians
to dimmer pixels where the MSE gradient is weaker.
Competitive pressure from overlapping Gaussians — when several are seeded on
the same bright stroke — drives rapid spatial diversification in the first
${\sim}500$ epochs via gradient-driven competition, achieving ${\approx}6\,\text{dB}$
higher PSNR by epoch $1{,}000$ than without-replacement seeding.

\paragraph{Coordinate inversion convention.}
Due to the affine-grid coordinate inversion of the renderer
(Section~\ref{sec:coord_convention}), the initialisation must set
$\Wraw_{k,5} = \atanh(-x_\text{pixel})$ and
$\Wraw_{k,6} = \atanh(-y_\text{pixel})$
so that the Gaussian renders at the desired pixel position.
This was the source of the coordinate-inversion bug (Section~\ref{sec:bug}).

\paragraph{Shape initialisation.}
All $K$ Gaussians are initialised with equal, small radii:
$\Wraw_{k,0} = \Wraw_{k,1} = -1.5$ (i.e.\ $\sigma_x = \sigma_y \approx 0.18$),
zero correlation ($\Wraw_{k,2} = 0$), and colour seeded from the pixel value
$c_k = \sigmoid(\logit(\I_{y_k, x_k}))$.

\paragraph{All-black images.}
If $\|\I\|_1 < 10^{-7}$ (all-black image), pixel-weighted sampling degenerates.
In this case positions are sampled uniformly, which is the correct behaviour.

\subsection{Optimisation}
\label{sec:optimisation}

We optimise $\Wraw$ with Adam and a cosine annealing schedule with warm restarts:

\begin{itemize}[leftmargin=2em]
  \item \textbf{Optimiser}: Adam, $\eta = 5 \times 10^{-3}$.
  \item \textbf{Scheduler}: CosineAnnealingWarmRestarts with $T_0 = \lfloor
        \text{epochs}/3 \rfloor$, $T_\text{mult} = 1$, $\eta_\text{min} = 10^{-5}$.
        Three full cosine cycles allow the learning rate to reset to its full value,
        helping the optimiser escape shallow local minima.
        Empirically, images that plateau at ${\sim}26\,\text{dB}$ can jump $6\text{--}9\,\text{dB}$
        when the LR resets to ${\approx}10^{-3}$.
  \item \textbf{Early stopping}: training halts when MSE\,$<10^{-4}$
        ($\PSNR \approx 40\,\text{dB}$).
  \item \textbf{Maximum epochs}: $3{,}000$.
\end{itemize}

Algorithm~\ref{alg:encode} gives the complete training procedure.

\begin{algorithm}[ht]
\caption{Gaussian encoder for a single image}
\label{alg:encode}
\begin{algorithmic}[1]
\Require Image $\I \in [0,1]^{H \times W}$, $K$, \textit{epochs}, $\eta$,
         \textit{early\_stop}, \textit{recycle\_every}
\State $\Wraw \leftarrow \textsc{InitGaussians}(\I, K)$
       \Comment{brightness-weighted, coord-inversion aware}
\State $\theta \leftarrow \textsc{Adam}([\Wraw],\; \eta)$
\State $\textit{scheduler} \leftarrow \textsc{CosineWarmRestarts}(\theta,\; T_0=\lfloor\textit{epochs}/3\rfloor)$
\For{$e = 0, 1, \ldots, \textit{epochs}-1$}
  \State $\W \leftarrow \textsc{ToPhysical}(\Wraw)$
  \State $\hat{\I} \leftarrow \Render(\W)$
         \Comment{differentiable Gaussian splatting}
  \State $\ell \leftarrow \MSE(\hat{\I}, \I)$
  \State Back-propagate $\ell$ and update $\Wraw$ via $\theta$
  \State Step scheduler
  \If{$(e+1) \bmod \textit{recycle\_every} = 0$}
    \State $\textsc{Recycle}(\Wraw, \I, \theta)$
           \Comment{teleport dead Gaussians to under-reconstructed regions}
  \EndIf
  \If{$\ell < \textit{early\_stop}$} \textbf{break} \EndIf
\EndFor
\State \Return $\textsc{ToPhysical}(\Wraw)$
\end{algorithmic}
\end{algorithm}

\subsection{Dead-Gaussian Detection and Recycling}
\label{sec:recycling}

A Gaussian is \emph{dead} when its centre lies on a background pixel
(target intensity $< 0.05$).  These Gaussians contribute nothing to the
reconstruction and waste optimiser capacity.

Every \textit{recycle\_every} epochs (default: $300$), we detect dead Gaussians
and \emph{teleport} them to under-reconstructed regions:

\begin{enumerate}[leftmargin=2em]
  \item Render the current approximation $\hat{\I}$.
  \item Compute residual $r = (\I - \hat{\I})_+ = \max(0,\;\I - \hat{\I})$.
  \item Sample new positions weighted by $r$ (regions the model is missing most).
  \item Re-initialise dead rows of $\Wraw$ with new positions and pixel colours.
  \item Reset the Adam moment tensors for recycled rows (stale second moments
        from the dead region would otherwise cause a large spurious update).
\end{enumerate}

\noindent
\textbf{Why not use colour threshold for death detection?}\ \
Overlapping Gaussians legitimately suppress each other's colour through
gradient-driven competition, so colour~$<0.05$ is a false positive.
Only the spatial criterion (centre on background) is used.

\subsection{The Vestigial \texorpdfstring{$\alpha$}{alpha} Parameter}
\label{sec:alpha}

The $\alpha$ (opacity) parameter is decoded from the raw parameters by
$\alpha = \sigmoid(\Wraw_{:,3})$, but it is \emph{never passed to the renderer}.
The function \texttt{generate\_2D\_gaussian\_splatting} does not accept an opacity
argument; Gaussian contributions are composited additively weighted only by $c$
(colour).
As a result, $\alpha$ is free to take any value during training without affecting
the rendered image.
This is verified experimentally in Section~\ref{sec:alpha_exp}.
